\section{Projective counterexamples with analytic density}
\label{uni-sec:examples}

The last construction prescribes the density before choosing the background. A symmetric pair of radial perturbations will have the same determinant, equal to a fixed polynomial in the global projective coordinate $t$. Patching the scalar coefficients while preserving their product makes the density real analytic on all of projective space. Its two quadratic zero hypersurfaces retain the separation that allows different extremal values.

\begin{theorem}\label{uni-thm:analytic-example}
Put
\[
 t=\frac{|A|^2}{|A|^2+|B|^2},\qquad
 a=\frac1{1+\mathrm{e}},\qquad b=\frac{\mathrm{e}}{1+\mathrm{e}},\qquad
 f_*(t)=(t-a)^2(t-b)^2.
\]
There exist a
smooth positive nonclosed $\omega$ and distinct smooth normalized
$\omega$-plurisubharmonic functions $u,v$ with
\begin{equation}\label{uni-eq:analytic-example}
 (\omega+\ddc u)^3=(\omega+\ddc v)^3=f_*(t)\omega_X^3.
\end{equation}
Both solution forms are positive away from the two connected real analytic
hypersurfaces $\{t=a\}$, $\{t=b\}$, and have complex rank two on them.
The density is globally real analytic and has precisely quadratic
transverse zeros.  For every $n\geq4$ the same assertions hold on
$\mathbb{P}^3\times\mathbb{P}^{n-3}$, with the pulled-back density and the product
Fubini--Study volume, and with rank $n-1$ on the two product
hypersurfaces.
\end{theorem}

\begin{proof}
The rank-one formula \eqref{eq:numa-cube} reduces the construction to scalar coefficients with a prescribed product. By \eqref{eq:numa-fs},
\begin{equation}\label{eq:analytic-fs-cube}
 \omega_X^3=6(\sigma')^2P\wedge Q\wedge R,
 \qquad \sigma'=\sigma(r)\sigma(-r)>0.
\end{equation}

\subsection{Prescribing the determinant}
\label{subsec:analytic-prescribed-determinant}

On $U$, one has $t=\sigma(r)$. We first construct two positive triples with product $f_*(\sigma)(\sigma')^2$ on the middle band. Their horizontal factors will lose exactly one direction at each endpoint.

Put $d_*(r)=f_*(\sigma(r))(\sigma'(r))^2$.  Since
$a=\sigma(-1)$, $b=\sigma(1)$, define the positive smooth divided
differences
\[
 k_-(r)=\int_0^1\sigma'(-1+s(r+1))\,\mathrm{d}s,\qquad
 k_+(r)=\int_0^1\sigma'(1+s(r-1))\,\mathrm{d}s,
\]
and $H_*(r)=k_-(r)k_+(r)\sigma'(r)>0$.  The fundamental theorem of
calculus gives
\begin{equation}\label{uni-eq:prescribed-det}
 \sigma(r)-a=(r+1)k_-(r),\quad
 \sigma(r)-b=(r-1)k_+(r),\quad
 d_*=(1-r^2)^2H_*^2.
\end{equation}
These formulas supply the smooth positive extension at both endpoints,
without a division by a vanishing function.  The symmetry
$\sigma(-r)=1-\sigma(r)$ makes $d_*$, and hence $H_*$, even.

Choose the flat bump
\[
 p(r)=
 \begin{cases}\mathrm{e}^{-1/(1-r^2)},&|r|<1,\\0,&|r|\geq1,
 \end{cases}
 \qquad q=p'.
\]
All derivatives vanish at $r=\pm1$: differentiating gives a finite sum
of powers of $(1-r^2)^{-1}$ multiplied by the exponential, and the
exponential dominates every such power.  On $|r|<1$,
$q=-2rp/(1-r^2)^2$.  Define
\begin{align}
 z_b&=\sqrt{q^2+2H_*^2},\qquad
 S_* =\sqrt{r^2+4p^2+2z_b(1-r^2)^2},\label{uni-eq:baseline1}\\
 x_b&=(S_*-r)/2,\qquad y_b=(S_*+r)/2.\label{uni-eq:baseline2}
\end{align}
Both radicands are strictly positive everywhere, so these are single
smooth formulas, including at the endpoints.  They satisfy
\begin{equation}\label{uni-eq:baseline-identities}
 z_b>|q|,\qquad x_b-y_b=-r,\qquad
 x_by_b=p^2+\frac{z_b}2(1-r^2)^2.
\end{equation}
The coefficients $x_b,y_b$ are nonnegative, positive away from the
endpoints, with values $(1,0)$ at $-1$ and $(0,1)$ at $1$.
On the open band,
\[
 S_*^2-(2p+|r|)^2=2(1-r^2)^2(z_b-|q|)>0,
\]
so $x_b>p$, $y_b>p$.  Consequently both triples
\begin{equation}\label{uni-eq:twotriples}
 (x_b-p,y_b+p,z_b-q),\qquad (x_b+p,y_b-p,z_b+q)
\end{equation}
are strictly positive there.

Here is the determinant cancellation.  Set
$D_*=x_by_b-p^2$, $E_*=p(x_b-y_b)=-rp$.
Equations \eqref{uni-eq:baseline-identities} and the formula for $q$ give
\[
 D_*=(z_b/2)(1-r^2)^2,\qquad qD_*=z_bE_*.
\]
The two products in \eqref{uni-eq:twotriples} therefore equal
\begin{equation}\label{uni-eq:det-cancellation}
 (D_*+E_*)(z_b-q)=(D_*-E_*)(z_b+q)
       =\frac{D_*}{z_b}(z_b^2-q^2)=d_*.
\end{equation}
For $|r|\geq1$, $p=q=0$ and the same identity reads
$x_by_bz_b=d_*$.  Both perturbed forms have exactly the prescribed
determinant, and rank two at the endpoints.

\subsection{Extension over the projective ends}
\label{subsec:analytic-global-patch}

The coefficients just constructed have the right determinant but need not extend over the projective ends. We interpolate the two horizontal coefficients to those of $f_*(t)^{1/3}\omega_X$ away from the zero levels, and determine the radial coefficient by the prescribed determinant.

Choose a smooth even cutoff $0\leq\eta\leq1$, equal to one on
$|r|\leq3/2$ and zero on $|r|\geq2$.  On $|r|<3/2$ set
$(x,y,z)=(x_b,y_b,z_b)$.  On $|r|>5/4$, where
$f_*(\sigma(r))>0$, put
\begin{align*}
 \kappa&=f_*(\sigma(r))^{1/3},\\
 x&=\eta x_b+(1-\eta)\kappa\sigma(r),\\
 y&=\eta y_b+(1-\eta)\kappa\sigma(-r),\qquad z=d_*/(xy).
\end{align*}
These coefficients are smooth and positive on the outer region.  On the
overlap $5/4<|r|<3/2$, $\eta=1$, $p=q=0$, and
$x_by_bz_b=d_*$; hence the two definitions agree identically.  For
$|r|\geq2$ they equal
$\kappa(\sigma(r),\sigma(-r),\sigma'(r))$.  Thus
\[
 \omega_0=xP+yQ+zR
\]
extends smoothly across both omitted projective lines, being equal there
to $f_*(t)^{1/3}\omega_X$.  This is smooth positive near those lines since
$f_*(0)=f_*(1)=a^2b^2>0$.  Globally $\omega_0$ is semipositive,
positive away from $r=\pm1$, and rank two on these levels.

The function
\[
 F(r)=2\int_0^r p(s)\,\mathrm{d}s
\]
is constant on each tail and extends to a smooth nonconstant function on
$X$.  Its half-Hessian is
$H_0=\ddc F/2=p(P-Q)+qR$.  By
\eqref{uni-eq:det-cancellation}, the forms
$A_0=\omega_0-H_0$, $B_0=\omega_0+H_0$ satisfy
\[
 A_0^3=B_0^3=f_*(t)\omega_X^3
\]
on the open block set and hence globally.  Their ranks and positivity
are those already checked in \eqref{uni-eq:twotriples}.  The construction is
symmetric: $x(-r)=y(r)$ and $z$ is even.

\subsection{A positive equation background}
\label{subsec:analytic-repair}

As in the first example, an exact correction makes the average form strictly positive without changing the two solution forms. Here the radial coefficient is already positive; the correction restores the missing horizontal direction at each endpoint. Choose a parameter $\varepsilon>0$ as restricted below.

Let $h(r)=(r^2-1)\eta(r)$, extend it by zero on the tails, and set
$\rho=\varepsilon h$.  The coefficients of
$\omega=\omega_0+\ddc\rho$ are
\begin{equation}\label{uni-eq:repair}
 (x+\varepsilon h',\ y-\varepsilon h',\ z+\varepsilon h'').
\end{equation}
On $|r|\leq3/2$, $h'=2r$, $h''=2$.
For $0\leq r\leq3/2$, $x+2\varepsilon r>0$: the only zero of
$x$ is at $r=1$, where the correction is positive.
On $[-3/2,0]$, $x$ has a positive minimum $m_0$, so positivity
holds if $3\varepsilon<m_0$.  Reflection proves positivity of the
second coefficient, and the third is $z+2\varepsilon>0$.
On the compact transition region $3/2\leq|r|\leq2$, all original
coefficients have a common positive lower bound $m_1$.  If
$M_1$ bounds $|h'|+|h''|$, taking also
$\varepsilon<m_1/(M_1+1)$ preserves their positivity there.
For $|r|\geq2$, the background equals
$f_*(t)^{1/3}\omega_X>0$.  This proves global strict positivity.

The correction is exact, so $\mathrm{d}\omega=\mathrm{d}\omega_0$.  At $r=0$,
evenness of $S_*$ gives $x'(0)=-1/2$, $y'(0)=1/2$.  Hence
\eqref{eq:numa-differential} gives
\[
 \mathrm{d}\omega_0=(z_b(0)+1/2)\,\mathrm{d}r\wedge(Q-P)\ne0.
\]
Indeed, if $\mathrm{d}r\wedge(Q-P)=0$, wedging with $P$ and $\dc r$
would contradict $P\wedge Q\wedge R>0$.  Thus the positive background
is nonclosed.

Now put $\widetilde u=-F/2-\rho$,
$\widetilde v=F/2-\rho$.  Then
$\omega+\ddc\widetilde u=A_0$,
$\omega+\ddc\widetilde v=B_0$.
Subtract their respective suprema to obtain $u,v$ with common
normalization zero.  Their difference remains nonconstant because $F$
is nonconstant.

\subsection{Zeros and higher dimensions}

It remains to check the global regularity of the prescribed density and the geometry of its zero levels. Taking a product with projective space then preserves the determinant equality and supplies examples in every higher dimension.

The function $t$ is real analytic in every projective chart, being a
quotient of real polynomials with positive denominator.  Hence $f_*(t)$
is globally real analytic.  Its two zero levels are regular because
$t=\sigma(r)$, $\sigma'>0$, and $\mathrm{d}r\ne0$.
Each level is connected: it is the quotient of the product of two
three-spheres of fixed nonzero radii by simultaneous circle multiplication.
It is compact, and a hypersurface has zero smooth volume.  The divided
differences in \eqref{uni-eq:prescribed-det} show near the two levels that
\begin{equation}\label{uni-eq:quadratic-zero}
 f_*(\sigma(r))=(r+1)^2g_-(r),\qquad
 f_*(\sigma(r))=(r-1)^2g_+(r),
\end{equation}
respectively, with smooth strictly positive transverse factors.

For $n\geq4$, let $\theta$ be the Fubini--Study form on
$Y=\mathbb{P}^{n-3}$, let $\pi_1:X\times Y\to X$ and $\pi_2:X\times Y\to Y$ be the projections, and set on $X\times Y$
\[
 \omega_n=\pi_1^*\omega+\pi_2^*\theta,\qquad
 \Omega_n=\pi_1^*\omega_X+\pi_2^*\theta,\qquad
 u_n=\pi_1^*u,\quad v_n=\pi_1^*v.
\]
Only the term with three factors from $X$ survives in an $n$-fold
product.  Consequently
\[
 (\omega_n+\ddc u_n)^n
 =\binom n3\pi_1^*(A_0^3)\wedge\pi_2^*(\theta^{n-3})
 =f_*(\pi_1^*t)\Omega_n^n,
\]
and the same holds for $v_n$.  The solution ranks on the two product
levels are $2+(n-3)=n-1$.  Restriction to a slice recovers
$\mathrm{d}\omega\ne0$, so $\omega_n$ is nonclosed.  Normalization,
distinctness, and real analyticity are preserved by pullback.
\end{proof}

Write $\operatorname{Vol}$ for volume with respect to $\omega_X^3$, or $\Omega_n^n$ in dimension $n$. The quadratic zeros complete the sharpness assertion of
Theorem~\ref{uni-thm:threshold}.  In finitely many tubular charts, the density
is bounded above and below by positive multiples of the square of a real
transverse coordinate, while the smooth volume has a positive bounded
Jacobian.  Away from these charts the density is bounded below.  Thus, with $s>0$ a sublevel parameter and $\alpha>0$ an integrability exponent,
\begin{equation}\label{uni-eq:sharp-asymptotic}
 \operatorname{Vol}\{f_*\leq s\}\asymp\sqrt s\quad(s\downarrow0),
 \qquad f_*^{-\alpha}\in L^1\quad\Longleftrightarrow\quad\alpha<\tfrac12
 \quad(\alpha>0).
\end{equation}
Thus every weaker inverse-power hypothesis with positive exponent below
$1/2$ admits smooth nonunique examples in every $n\geq3$.

\begin{remark}
For this example, $U=\{f_*(t)>0\}$ has three components: the middle band and the two outer regions. The solution forms have the same missing direction on each zero level, so $E_{\rm unif}=U$. The maximum and minimum of their potential difference occur on the two outer components, whose closures are disjoint, in accordance with \cref{uni-thm:anchors}. Although the middle component meets both outer closures, this adjacency does not propagate an extremal value through it.
\end{remark}

